\subsection{Đánh giá kết quả thực nghiệm}

Trên cơ sở kết quả thực nghiệm trình bày ở mục~4.2, hệ thống được đánh giá tổng thể theo bốn tiêu chí: \textit{tính an toàn}, \textit{tính dễ sử dụng}, \textit{khả năng triển khai} và \textit{độ tin cậy và bảo mật}. Bốn tiêu chí này phản ánh đầy đủ các góc nhìn đối với một hệ thống bầu cử trực tuyến -- từ yêu cầu kỹ thuật của đội vận hành, trải nghiệm thực tế của cử tri đến yêu cầu kiểm chứng của bên giám sát độc lập.

\begin{itemize}
    \item \textbf{Tính an toàn:} Hệ thống đáp ứng đồng thời sáu tính chất của một cuộc bầu cử an toàn -- \textit{bí mật}, \textit{duy nhất}, \textit{xác thực}, \textit{không thể chối bỏ}, \textit{kiểm chứng được} và \textit{ẩn danh} -- thông qua sự phối hợp của nhiều tầng phòng vệ độc lập. Lựa chọn của cử tri được ký mù theo lược đồ EC-Schnorr trước khi gửi đi nên máy chủ không quan sát được nội dung phiếu; phiếu được lưu lên sổ cái Fabric dưới dạng \textit{commitment} trong cây Merkle, cho phép kiểm chứng tính toàn vẹn mà không lộ ánh xạ tới cử tri. Quản lý nonce trên Redis loại trừ bỏ phiếu trùng và tấn công phục hồi khóa qua nonce reuse, còn JWT hai cặp khóa ngăn giả mạo định danh xuyên suốt phiên làm việc.

    \item \textbf{Tính dễ sử dụng:} Cả hai ứng dụng được thiết kế tối thiểu hóa số thao tác trên mỗi màn hình. Phía cử tri, quy trình bỏ phiếu rút gọn còn bốn bước tuần tự (đăng nhập, chọn cuộc bầu cử, chọn lựa chọn, xác nhận); biên lai và kết quả được lưu trữ ngay trong ứng dụng, không yêu cầu cử tri ghi nhớ thêm thông tin kỹ thuật. Phía quản trị viên, giao diện dạng bảng điều khiển và phân chia tab theo từng vai trò chức năng giúp luồng vận hành rõ ràng. Các thao tác mật mã được thực hiện ngầm phía nền tảng, người dùng không cần kiến thức về chữ ký mù hay cây Merkle để hoàn thành phiếu bầu hoặc xác minh.

    \item \textbf{Khả năng triển khai:} Hạ tầng blockchain được điều phối qua ChainLaunch, cho phép dựng mạng Hyperledger Fabric đầy đủ trong vài phút thay vì hàng giờ cấu hình thủ công; cơ chế Infrastructure as Code qua Terraform giúp tái lập môi trường thử nghiệm và sản xuất nhất quán. Backend NestJS đơn khối theo mô đun cho phép tách dần thành dịch vụ vi mô khi nhu cầu mở rộng tăng. Ứng dụng di động trên Expo từ một mã nguồn duy nhất xuất bản đồng thời cho Android và iOS, giảm chi phí phát triển và bảo trì.

    \item \textbf{Độ tin cậy và bảo mật:} Độ tin cậy được bảo đảm bởi đặc tính bất biến của sổ cái Fabric kết hợp đồng thuận Raft trong môi trường nhiều tổ chức ký, loại trừ khả năng đơn phương sửa đổi dữ liệu bầu cử. Cây Merkle gắn với cam kết bất biến cho phép cử tri và bên kiểm toán tự xác minh sự tồn tại của phiếu mà không cần tin tưởng máy chủ vận hành. Khóa bí mật của các nút ký được mã hóa AES-256-GCM ở tầng lưu trữ tĩnh; HTTP đối ngoại được cứng hóa bằng Helmet, CORS, Rate Limiting và bắt buộc HTTPS; giao tiếp nội bộ giữa backend và các nút ký được bảo vệ bằng mTLS với chứng chỉ ngắn hạn.
\end{itemize}

Nhìn chung, kết quả thực nghiệm cho thấy hệ thống đã hiện thực hóa được các mục tiêu thiết kế đặt ra ở Chương~2 và các tầng bảo mật trình bày ở Chương~3: vừa giữ tính ẩn danh và bí mật của phiếu bầu, vừa duy trì khả năng kiểm chứng độc lập cho mỗi cử tri, đồng thời vẫn đảm bảo trải nghiệm sử dụng đơn giản và chi phí triển khai phù hợp với các tổ chức cỡ trung. Đây là cơ sở khẳng định tính khả thi của hướng tiếp cận kết hợp blockchain với mật mã chữ ký mù tập thể trong các bài toán bỏ phiếu thực tế.
